% భాగం: set-theory
% అధ్యాయం: ordinals
% విభాగం: introduction

\documentclass[../../../include/open-logic-section]{subfiles}

\begin{document}

\olfileid{sth}{ordinals}{intro}

\olsection{పరిచయం}

\olref[z][]{chap}లో,
\stagesinf{} రూపంలో
సమితుల స్థాయిక్రమానికి
అనంతవ దశ ఉందని
ప్రతిపాదించాం (మన
అనంతత్వ స్వీకృతాన్ని
కూడా చూడండి). అయితే
\stagessucc{} ఉంటే,
అనంతవ దశ వద్ద
ఆగలేం; ముందుకు
సాగాల్సిందే. కాబట్టి
తొలి అనంత దశ తరువాతి
దశలో, తొలి అనంత
దశలో అందుబాటులో ఉన్న
సమితులతో ఏర్పడగల
అన్ని సమాహారాలను
ఏర్పరుస్తాం; మళ్లీ;
మళ్లీ; మళ్లీ; \dots

ఇక్కడ మనం అంతర్లీనంగా
ఒక ``సహజ'' సంఖ్య
భావాన్ని వాడడం
ప్రారంభించాం: అన్ని
సహజ సంఖ్యల \emph{తరువాత}
కూడా సంఖ్యలు ఉండవచ్చనే
భావాన్ని. ప్రత్యేకంగా
అది \emph{అతిపరిమిత
క్రమసంఖ్య} భావం.
ఈ భావాన్ని మరింత
కచ్చితంగా చెప్పడమే
ఈ అధ్యాయం లక్ష్యం.
క్రమసంఖ్య అనే సాధారణ
భావాన్ని పరిశీలించి,
కొన్ని సమితులను మన
క్రమసంఖ్యలుగా స్పష్టంగా
నిర్వచిస్తాం.

\end{document}
